% =====================================================================
%  TRIADE — édition académique v0.13.0 (manuscrit IMRAD, anglais)
%  Corps du document : compilation Tectonic (XeTeX).
%  Les tableaux de données sont GÉNÉRÉS par scripts/extraire_donnees_v130.py
%  (donnee unique : docs/JOURNAL_TRAVAIL.md + docs/DECISIONS.md + les
%  sorties brutes des instruments preuves_v110/v120/v130) et \input
%  ci-dessous — aucune donnée n'est tapée à la main ; les comptes du journal
%  vivent en macros (donnees/macros_v130.tex, régénérées à chaque passe).
% =====================================================================
\documentclass[11pt,a4paper]{article}

\usepackage{fontspec}
\setmainfont{Liberation Serif}
\setsansfont{Liberation Sans}
\setmonofont{Liberation Mono}
\usepackage[french,english]{babel}

\usepackage{amsmath,amssymb}
\usepackage{geometry}
\geometry{a4paper, top=2.6cm, bottom=2.4cm, left=2.3cm, right=2.3cm}
\usepackage{graphicx}
\usepackage{xcolor}
\definecolor{liaison}{HTML}{3681A6}
\usepackage{booktabs}
\usepackage{tabularx}
\newcolumntype{Y}{>{\raggedright\arraybackslash}X}
\usepackage{enumitem}
\setlist{noitemsep, topsep=2pt, leftmargin=1.4em}
\usepackage{microtype}
\usepackage{url}
\usepackage[numbers,sort&compress]{natbib}
\bibliographystyle{unsrtnat}
\setlength{\bibsep}{2pt}
\usepackage{fancyhdr}
\usepackage[colorlinks=true, linkcolor=liaison, citecolor=liaison,
            urlcolor=liaison, unicode=true,
            pdftitle={TRIADE: Enabling LLM Agents to Write Native Code by Hand Under Signed-Intent Governance (Academic Edition v0.13.0)},
            pdfauthor={Alexis Lopez}]{hyperref}
\pagestyle{fancy}
\fancyhf{}
\fancyhead[L]{\small\itshape TRIADE — academic edition v0.13.0}
\fancyhead[R]{\small\thepage}
\renewcommand{\headrulewidth}{0.3pt}

\newcommand{\code}[1]{\texttt{\small #1}}
\newcommand{\RQ}[1]{\textbf{RQ#1}}

% Les comptes EXACTS du journal (régénérés par extraire_donnees_v130.py —
% jamais un arrondi, jamais une divergence abstract/intro ; D190)
% Fichier GÉNÉRÉ par extraire_donnees_v130.py — ne pas éditer à la main.
% Les comptes EXACTS, extraits à l'instant (jamais un arrondi).
\newcommand{\nbLignesJournal}{9{,}417}
\newcommand{\nbLignesDecisions}{4{,}539}
\newcommand{\plageDecisions}{D1--D192}
\newcommand{\nbEntreesJournal}{205}


\title{\textbf{TRIADE: Enabling LLM Agents to Write Native Code by Hand\\
Under Signed-Intent Governance\\
\large A longitudinal single-case study over eleven milestones, twelve field
runs, and one sealed application}}
\author{Alexis Lopez\thanks{Independent researcher; corresponding author.
Project repository: \code{alexisfrlopez-debug/triade} (GitHub mirror).}}
\date{11 November 2026 — Academic edition v0.13.0 (corrected statistics, audited benchmark, expanded related work)}

\begin{document}
\maketitle
\thispagestyle{fancy}

% ------------------------------------------------------------------
\begin{abstract}
\noindent Large language model (LLM) agents currently produce native code by
generating high-level source and delegating to a compiler toolchain; the
translation from intent to machine code passes through layers the agent
neither writes nor verifies. We report on TRIADE, a development protocol in
which an LLM agent writes native x86-64, RISC-V, and TLS~1.3 code
\emph{directly as assembly notation}, governed by signed layer-I
specifications, an encrypted and fragmented card that binds intentions to
byte ranges, and an epistemology that separates \emph{oracle} evidence
(differential emulation in the laboratory) from \emph{terrain} evidence
(execution on the developer's real machine). We present a longitudinal
single-case study of the protocol's validation programme: eleven milestones
(M0--M10, 26 October -- 11 November 2026), culminating in ``Loupe'', a real
Windows hex-editor application written entirely under the protocol and
delivered to its user. We quantify four research questions on data drawn
from a \nbLignesJournal-line append-only engineering journal: (RQ1) the
escaped-defect rate of changes that had passed all laboratory controls ---
eleven of twelve field iterations surfaced at least one escaped defect,
each named to a class and a mechanism before any fix; (RQ2) the measured
token cost per validated change --- 51 tokens to locate a block, 719 to
edit one, 1{,}088 to add one, against design targets of 64/3{,}000/6{,}000;
(RQ3) the defect classes structurally invisible to emulation, which we
enumerate with their evidence (loader and PE-format semantics, ABI
shadow-space and SSE alignment, Windows error-code conventions,
harness-embedded assumptions); and (RQ4) the safety of hot-patching sealed
code --- zero post-patch regression across the milestone campaigns, with
signed version monotonicity and an armed anti-rollback check. The study is
single-case (N=1), its oracle is derived from the same specifications it
checks, and no baseline compiler-based replication exists yet; we state
these limits explicitly. We claim as contributions the \emph{integration}
of contract-first notation, signed-manifest governance and oracle/terrain
verification into one pipeline for agent-first development --- not the
underlying cryptographic or testing primitives, which are standard.
\end{abstract}

\begin{otherlanguage}{french}
\begin{abstract}
\noindent Les agents de modèles de langage produisent aujourd'hui du code
natif en générant du source de haut niveau délégué à une chaîne de
compilation ; la traduction de l'intention vers le code machine traverse
des couches que l'agent n'écrit ni ne vérifie. Nous rapportons TRIADE, un
protocole de développement dans lequel un agent écrit du code natif
x86-64, RISC-V et TLS~1.3 \emph{directement en notation d'assemblage},
gouverné par des spécifications signées (couche I), une carte chiffrée et
fragmentée qui relie les intentions aux plages d'octets, et une
épistémologie qui sépare la preuve d'\emph{oracle} (émulation
différentielle au laboratoire) de la preuve de \emph{terrain} (exécution
sur la machine réelle du développeur). Nous présentons une étude de cas
unique longitudinale du programme de validation : onze jalons (M0--M10, du
26 octobre au 11 novembre 2026), culminant avec « Loupe », une véritable
application Windows (éditeur hexadécimal) écrite entièrement sous le
protocole et livrée à son utilisateur. Nous quantifions quatre questions
de recherche sur des données tirées d'un journal d'ingénierie append-only
de \nbLignesJournal{} lignes : (RQ1) le taux de défauts échappés des changements ayant
passé tous les contrôles de laboratoire — onze itérations de terrain sur
douze ont révélé au moins un défaut échappé, chacun nommé par classe et
mécanisme avant toute correction ; (RQ2) le coût mesuré en tokens par
changement validé — 51 tokens pour localiser un bloc, 719 pour l'éditer,
1~088 pour l'ajouter, contre des cibles de conception de 64/3~000/6~000 ;
(RQ3) les classes de défauts structurellement invisibles à l'émulation,
énumérées avec leurs preuves (sémantique du chargeur et du format PE,
espace d'ombre et alignement SSE de l'ABI, conventions de codes d'erreur
Windows, hypothèses incrustées dans le harnais) ; et (RQ4) la sûreté du
patch à chaud de code scellé — zéro régression post-patch sur les
campagnes de jalons, avec monotonie de version signée et anti-rollback
armé. L'étude est un cas unique (N=1), son oracle dérive des mêmes
spécifications qu'il vérifie, et aucune réplication avec compilateur de
référence n'existe encore ; ces limites sont énoncées explicitement. Nous
revendiquons comme contribution l'\emph{intégration} — notation à
contrats d'abord, gouvernance à manifestes signés, vérification
oracle/terrain — en un seul pipeline pour le développement piloté par
agent, et non les primitives cryptographiques ou de test sous-jacentes,
qui sont standards.
\end{abstract}
\end{otherlanguage}

% ------------------------------------------------------------------
\section{Introduction}\label{sec:intro}

An LLM-based agent that develops software today almost always works at the
level of high-level source code, and the mapping from that source to
machine code is delegated to a compiler toolchain. This division of labour
has an epistemological consequence: the agent's \emph{intent} is checked
against source text, while the \emph{bytes that execute} are produced by a
translation chain the agent does not write, does not see, and cannot
re-derive. Bugs in that mapping, and drift between what was specified and
what was produced, become a class of defects that no amount of source-level
review can exclude.

TRIADE (\emph{Triple Mirror: Intention, Bytes, Card}) is a protocol
designed for the opposite position: the agent writes native code
\emph{directly}, as assembly notation assembled by a deterministic,
verifiable assembler path --- there is no optimizing or opaque translation
chain; what the agent writes is, by construction, what executes
(assembly notation is verified by an exact round-trip against the emitted
bytes at every build). The protocol was developed and validated by a
single human developer directing an LLM agent, over eleven milestones
between 26 October and 11 November 2026, and is documented in a
\nbLignesJournal-line append-only work journal, a \nbLignesDecisions-line
append-only decision log (\plageDecisions), and sealed artifacts whose
SHA-256 fingerprints are consigned at each sealing ceremony.

This paper is the first academic edition of that effort. Its purpose is
not to argue that hand-written machine code is superior to compiled code
in general; it is to report, with the data the protocol itself produced,
what happens when an LLM agent develops real native software this way ---
what it costs, what breaks, what laboratory verification could not see,
and what governance was needed to keep the process falsifiable. The
protocol's cryptographic and testing primitives (signature schemes,
secret sharing, differential testing, key ceremonies) are standard and
cited; the contribution claimed is their \emph{integration} into one
agent-first pipeline, together with the oracle/terrain epistemology and
the signed-intent governance described below.

The study is structured around four research questions:

\begin{itemize}
\item \RQ{1} \emph{Escaped-defect rate.} When a change passes every
laboratory control (differential oracles, harnesses at three fixed seeds,
structural checks, seal guards), how often does it still fail on the real
target machine, and at what granularity can the failure be named?
\item \RQ{2} \emph{Token cost per validated change.} What does it cost,
in LLM context tokens, to locate, edit, or add a block of sealed native
code under the protocol's card-based indexing, measured rather than
estimated?
\item \RQ{3} \emph{Emulation-invisible defect classes.} Which classes of
defects are structurally invisible to a Unicorn-based emulation harness,
and therefore escape to the terrain regardless of laboratory coverage?
Framed differently: a compiler silently guarantees several of these
classes --- RQ3 quantifies what must be re-derived, and learned, when the
compiler is removed from the agent's path.
\item \RQ{4} \emph{Hot-patch safety.} Can code be patched in place while
sealed artifacts remain byte-stable, with zero post-patch regression and
an armed rollback barrier, across ISAs?
\end{itemize}

The remainder of the paper describes the protocol and the study methods
(Section~\ref{sec:methods}), the results per research question
(Section~\ref{sec:results}), a discussion (Section~\ref{sec:discussion}),
threats to validity (Section~\ref{sec:threats}), related work
(Section~\ref{sec:related}), and declarations.

% ------------------------------------------------------------------
\section{Methods}\label{sec:methods}

\subsection{The TRIADE protocol in brief}\label{sec:protocole}

TRIADE organises development in three indissociable layers. The
\emph{intention layer} (layer I) holds one normalized specification per
code block: identifier, intent, pre/post-conditions and invariants, named
failure cases, verification vectors, fixtures, provenance and the
decisions that shaped it. The \emph{block layer} (layer B) holds the
assembly notation of each block, at fixed addresses, with deterministic
assembly and a mandatory round-trip check (assemble, disassemble,
re-assemble; any divergence refuses the build). The \emph{card layer}
(layer C) is an encrypted, fragmented, signed index (the \emph{card}) that
binds every block to its byte range, contract and fingerprints; only the
card is needed to locate or edit a block, never the whole binary.

Eleven rules (R1--R11) fix the discipline; the load-bearing ones for this
study are: R2 (\emph{written before code}: every block is specified before
any notation byte exists); R4 (\emph{stable contract}: a change may not
silently alter a block's contract); R5 (\emph{verify then write}: the
agent must re-derive the check before writing); R6 (crystallisation: an
artifact may only be sealed after ten consecutive identical rebuilds);
R7 (\emph{sealed inheritance}: inherited bytes live at their sealed
addresses, byte-identical); R10 (load-time verification: checksums per
block, signed version monotonicity, anti-rollback floor); R11 (no
optimization without measurement).

The development loop itself (update an intention, re-derive the oracle,
edit the notation, rebuild, verify, seal) and the seal ceremony model are
described in the protocol specification v0.13.0, which this edition
accompanies; we summarise only what the evaluation needs.

\paragraph{Governance: signed intent and append-only records.}
Every layer-I tree is condensed and signed with an Ed25519 key kept
strictly offline; a build refuses an unsigned layer I. Structural
decisions are consigned in an append-only decision log --- each entry has
an identifier, date, status, context, the retained choice, rejected
alternatives and amendment links; entries are never rewritten or
back-dated. The fourteen signed trees at the time of the study are verified
by a guard that must pass \emph{before} any build (14/14 at the last
run). Sealing a release is a \emph{fact of the card}: the sealed tag
lives in the card's build fragment; the artifact's bytes never move
(Table~\ref{tab:empreintes} lists the sealed artifacts and their
fingerprints).

\paragraph{The oracle/terrain epistemology.}
The protocol distinguishes two evidence regimes. \emph{Oracle evidence} is
produced in the laboratory: differential oracles derived from the layer-I
specifications (never exported, always regenerated from I), scenario
oracles that replay expected screens byte-for-byte, structural oracles
that read the composed image instruction by instruction, and Unicorn-based
emulation harnesses with Windows ABI shims. \emph{Terrain evidence} is
produced on the developer's real machine: the sealed artifact, a
scenario-loader script, and witnesses (byte-exact logs, showcase runs)
that are pushed back to an evidence vault and cross-checked (``recoup'')
against the laboratory's expectations. The study's unit of failure is the
\emph{named defect}: a failure is not counted until a class, a mechanism
and an evidence piece have been assigned to it, and every fix is a
documented surgical revision (kit-only or piece-only, with byte diffs).

\paragraph{The agent.}
The developer directed an LLM agent (Z.ai GLM family, ``the laboratory'')
through short natural-language instructions; the agent executed all
laboratory work (notation, oracles, harnesses, builds, ceremonies,
consignation) in a sandboxed Linux environment with no access to the
offline signing key. The complete transcript of the collaboration is the
append-only journal and work log themselves --- every session ends with a
consigned entry listing what was done, what was verified and what remains
--- plus the decision log; these, with the evidence vault, are the
primary sources of this paper. The agent is described here and in the
acknowledgements; it is not an author.

\subsection{The validation programme and the case study}\label{sec:cas}

The protocol was validated by eleven milestones, each with measured exit
criteria (Table~\ref{tab:jalons}); M0 established the artifact chain,
M2 hot patching on a real machine, M3 a million-case fuzzing campaign,
M4 an encrypted Schannel tunnel, M6/M7 portability to Linux and RISC-V,
M8 direct syscalls, M9 pure TLS~1.3 (28 blocks, official vectors from
FIPS~180-4, RFC~2104/4231, RFC~5869, RFC~7748, SP~800-38D, handshake
replayed on the RFC~8448 transcript, and a real encrypted session with a
real server), and M10 the full-scale test: \emph{Loupe}, a real Windows
hex-editor application --- 18 blocks in five families (file, digests,
rendering, navigation, editing) --- written entirely under the protocol,
delivered as a field kit to the developer's machine, and run twelve times
(Table~\ref{tab:runs}). Verification volumes per milestone are
consolidated in Table~\ref{tab:oracles}.

\subsection{Measurements}\label{sec:mesures}

\emph{Token budgets} were instrumented at M1 (decision D31/D34) with the
\code{cl100k\_base} tokenizer over the sealed card and real sources:
tokens to locate a block (read one \code{frag-loc} entry), to edit one
(additional \code{frag-con} + notation), to add one (additional
profile thresholds). \emph{Escaped defects} are counted per terrain
iteration under an operational rule defined once and applied uniformly:
a run counts toward the escaped-defect rate if at least one defect that
had passed every laboratory control in force at the time was named (a
class, a mechanism, an archived evidence piece) before any fix ---
kit-side and piece-side defects both count, the deliverable being the
kit+piece ensemble; runs~6 and~7 each surfaced two such defects and each
counts \emph{once} in the run rate (the rate measures escape incidence
per delivery, not per defect); the run-7 [D54-QUATRI\`EME] recurrence
is counted as a manifestation of the deposit-drift class first named in
the M4-era campaign, not as a new class; and the run-12 [T0-P\'ERIM\'E]
residue --- a stale banner string touching no gate, disclosed in the run
table --- is \emph{not} counted: it violates no control's contract and
alters no behaviour any control guards, and counting a text refresh as a
functional escape would misstate the rate's meaning.
\emph{Hot-patch safety} is measured as post-patch regressions (target:
zero) and as the byte-stability of sealed artifacts across patch
campaigns. All numbers reported in Section~\ref{sec:results} are
generated programmatically from the journal and decision log by the
extraction script shipped with this edition; the manuscript contains no
hand-typed data value.

% --- Tableaux générés depuis le journal (source unique de vérité) ---
% Fichier GÉNÉRÉ par extraire_donnees_v130.py — ne pas éditer à la main.
% Les cinq tableaux historiques : extraits/validés par le module extraire_donnees_v100.py (ancres du journal, empreintes mesurées).
% La table des preuves : PARSÉE depuis preuves_v110/v120/v130 — JAMAIS une troncature (la dette de la Table 6, soldée).

\begin{table}[htbp]\centering\footnotesize
\caption{Milestones M0--M10 of the TRIADE validation programme (26 Oct -- 11 Nov 2026). Every milestone was closed on measured exit criteria; sources: work journal entries and decision log.}
\label{tab:jalons}
\begin{tabularx}{\textwidth}{@{}ll lX@{}}
\toprule
\textbf{Milestone} & \textbf{Sealed tag} & \textbf{Date} & \textbf{Lab evidence (exit criteria met)} \\
\midrule
\textbf{M0} & \code{0.1.0} & 10-26 & round-trip 100\\%; executable < 4 KiB \\
\textbf{M1} & \code{0.2.0} & 10-26 & oracle 21/21; first crystallisation \\
\textbf{M2} & \code{0.3.0} & 10-27 & conformity 68/68; 8 instrumented iterations \\
\textbf{M3} & \code{0.4.0} & 10-28 & $10^6$ fuzzed cases $\times$2 seeds; 161 conformity controls \\
\textbf{M4} & \code{0.5.0} & 10-29 & 560-case SSPI scenario oracle $\times$4 seeds; 12 mutations caught \\
\textbf{M5} & \code{0.6.0} & 11-04 & per-build fingerprints active \\
\textbf{M6} & \code{0.7.0} & 11-05 & 7/7 conformity suite; layer-I diff zero \\
\textbf{M7} & \code{0.8.0} & 11-05 & round-trip 100 \% on new ISA; fence.i patch \\
\textbf{M8} & \code{0.9.0} & 11-05 & 61 harness greens; 131 socket-socle controls \\
\textbf{M9} & \code{1.0.0} & 11-07 & 28 blocks; 589 structural locks; handshake 54/54 $\times$3 seeds; cross-lock 86/86 \\
\textbf{M10} & \code{loupe 0.1.0 / 0.1.1} & 11-11 & 18 blocks; 159 controls $\times$3 seeds; cross-lock 26/26; seal guard 14/14 \\
\bottomrule
\end{tabularx}\end{table}

\begin{table}[htbp]\centering\footnotesize
\caption{Measured token budgets (M1 instrumentation, tokenizer cl100k\_base) against design targets, and generation-1 lineage volumes (D34). The lineage TOTAL counts specifications and delivered notations only --- drafts are excluded ($8{,}774 + 5{,}099 = 13{,}873$).}
\label{tab:budgets}
\begin{tabular}{@{}lrrl@{}}
\toprule
\textbf{Operation} & \textbf{Target} & \textbf{Measured} & \textbf{Source} \\
\midrule
Locate a block & 64 & 51 & D31/D34, M1 measurement (cl100k\_base) \\
Edit a block & 3000 & 719 & D31/D34, M1 measurement \\
Add a block & 6000 & 1,088 & D34, M1 (draft 1086, +2 drift re-measured) \\
Whole lineage (specs) & --- & 8,774 & D34, generation-1 card \\
Whole lineage (drafts) & --- & 3,973 & D34 \\
Whole lineage (delivered notations) & --- & 5,099 & D34 \\
Whole lineage (TOTAL) & --- & 13,873 & D34 \\
\bottomrule
\end{tabular}\end{table}

\begin{table}[htbp]\centering\footnotesize
\caption{Verification instruments and volumes per milestone. All numbers are drawn from the corresponding journal entries.}
\label{tab:oracles}
\begin{tabularx}{\textwidth}{@{}l X@{}}
\toprule
\textbf{Milestone} & \textbf{Instrument (measured volume)} \\
\midrule
\textbf{M1} & mini-libc oracle — 21/21 \emph{(D32; tag 0.2.0)} \\
\textbf{M2} & echo conformity + differential — 68/68; 8 iterations, 8 defect classes \emph{(D43; tag 0.3.0)} \\
\textbf{M3} & fuzz campaign + profile conformity suite — 1,000,000 cases $\times$2 seeds; 161 controls \emph{(D72/D73)} \\
\textbf{M4} & SSPI scenario oracle + per-block differential — 560 cases $\times$4 seeds; 10,000 fuzzed cases $\times$2 seeds per block; 165k leaf differentials \emph{(D68-D93)} \\
\textbf{M5} & build fingerprint registry + verifier — per-build sha256 registers \emph{(D134, D148)} \\
\textbf{M6} & profile conformity suite (Linux) — 7/7; lab recoup 22/22 \emph{(D154)} \\
\textbf{M7} & RISC-V round-trip + cross-profile lock — 100 \% round-trip; 29/29 recoup \emph{(D155)} \\
\textbf{M8} & sys harness on synthetic ntdll fixtures — 61 greens; 131 socle controls; 30/30 recoup \emph{(D160-D165)} \\
\textbf{M9} & tls13 structural oracle + handshake harness + cross-lock — 589 structural locks; handshake 54/54 $\times$3 seeds; cross-lock 86/86; 560+365+131 replayed \emph{(D166-D170)} \\
\textbf{M10} & Loupe: 3 Unicorn harnesses + structural oracle + seal guard — 159 controls $\times$3 seeds (61+53+45); 121 structural locks; cross-lock 26/26; seal guard 14/14; crystallisation 26 x 10/10 \emph{(D172-D181)} \\
\bottomrule
\end{tabularx}\end{table}

\begin{table}[htbp]\centering\scriptsize
\caption{The twelve terrain runs of the Loupe field test (M10). Every failure was \emph{named} (a class, a mechanism, an evidence piece) before being fixed; every fix was documented as a surgical revision. The sealed artifacts never changed after their sealing run. Sources: journal entries 204--222, decisions D178/D181.}
\label{tab:runs}
\begin{tabularx}{\textwidth}{@{}r l X X@{}}
\toprule
\textbf{\#} & \textbf{Kit} & \textbf{Outcome (named defect)} & \textbf{Documented fix} \\
\midrule
1 & \code{dev.1} & loader refusal 0xc000007b. Missing BASERELOC: ASLR image without relocation table (lab harness parses, never loads) & dev.2: minimal 12-byte BASERELOC, 7 bytes in 6 ranges \\
2 & \code{dev.2} & same refusal 0xc000007b. Necessary-not-sufficient: page hole 0x3000-0x4000 + massive BSS outside the proven envelope of 8 prior artifacts & dev.3: image within envelope (VS==RS, no BSS) \\
3 & \code{dev.3} & loader mapped the image (decisive step); T0 witness failed. [CADRE-NE]: .data initializer wrong (sentinel at +8 - a SIZE read as an offset instead of +528) & dev.4: initializer = exact mirror of harness CADRE\_VIERGE; guard healed (control derives from plan, never from the piece it guards) \\
4 & \code{dev.4} & T0 all green; T1 named -6; showcase died instantly. Stack-argument class: msx64 args 5+ at home space instead of [rsp+0x20]+ (loupe-01/05/06) + SSE alignment (loupe-10/11) & dev.5: slots fixed; msx64 ABI audit born at build (static, catches SSE class); falsification control V7 added \\
5 & \code{dev.5} & first run acquitting the piece; T1 named -7 VERROUILLE. [SESSION]: kit itself - scenario-as-loader holds loupe.exe as a live section, Windows refuses gle 32 to open it for write & dev.6 KIT-ONLY (artifact unchanged): byte-copy shadow temoin-session.bin \\
6 & \code{dev.6} & two defects from one witness. False save error + missing showcase thread; [HINSTANCE-PEB]: hInstance read from PEB instead of GetModuleHandleW(NULL) & piece dev.6 (r12 = GetModuleHandleW(NULL)) + kit dev.7 \\
7 & \code{dev.7} & T0-T5 all green (save proven on terrain); then two defects. [EAX-SANS-SIGNE] at kit (unsigned 32-bit reads of negative codes) + [D54-QUATRIEME] at piece (4 contiguous FirstThunk ranges without NULL separators) & piece dev.7 (IAT relocated 0x3100, 4 ranges separated by NULL, 63 call sites repointed) + kit dev.8 \\
8 & \code{dev.8} & T0-T7 all green for the FIRST time (functional contract fully proven). Probe died on its own Python; showcase still mute (5th death) & kit dev.9 (probe made load-bearing) \\
9 & \code{dev.9} & T0-T7 all green second time; 6th showcase death named by system then by lab. [OFT-QWORD]: 4 OFT tables at QWORD stride instead of DWORD - the loader walks 10.18.6.3 entries at DWORD stride; first cause of ALL deaths found on the piece bytes & piece dev.8 (89 bytes in 52 ranges, sealed 0.1.0 candidate) + kit dev.10 \\
10 & \code{dev.11} & T0-T7 all green third time; SHOWCASE ALIVE FIRST TIME (0 and 0, no WER). [FREE-SANS-TYPE] at kit: FreeLibrary called without argtypes, 64-bit HMODULE truncated to C int - ArgumentError killed cascade 205 & kit dev.12 (argtypes=[c\_void\_p], except Exception widened, probe armored try/finally) \\
11 & \code{dev.12} & T0-T7 all green 4th time; cascade 205 complete 1st time; showcase alive 2nd time; green holds entire. [SONDE-EXE] named by lab on pieces: structural false negative - library loader NEVER resolves an EXE's imports; liaison proven by the living showcase & seal loupe 0.1.0 on recouped run 39/39 (D178): notation debt paid (587$\rightarrow$0), crystallisation 26x10/10, anti-rollback armed \\
12 & \code{dev.1} & T0-T7 all green 5th time; cascade 205 2nd time; showcase alive 3rd time (0 and 0, first of piece 0.1.1). [T0-PERIME] residual named (dev.12-era banner text - no gate touched) & seal loupe 0.1.1 (D181): [R10-PERIMEE] debt paid at birth (17 sums refreshed), card generation 17, artifact intact \\
\bottomrule
\end{tabularx}\end{table}

\begin{table}[htbp]\centering\footnotesize
\caption{Sealed artifacts (SHA-256, verified byte for byte against the card and journal at edition time).}
\label{tab:empreintes}
\begin{tabular}{@{}l r l@{}}
\toprule
\textbf{Artifact} & \textbf{Bytes} & \textbf{SHA-256 (first 16 hex)} \\
\midrule
\code{loupe.exe 0.1.1} & 127,488 & \code{c111510535767a40\ldots} \\
\code{loupe.exe 0.1.0} & 127,488 & \code{054de960c31425d0\ldots} \\
\code{proxy.exe 1.0.0} & 65,024 & \code{446279f3ab356900\ldots} \\
\bottomrule
\end{tabular}\end{table}

\begin{table}[htbp]\centering\scriptsize
\caption{Strengthened evidence across editions, every row generated from the raw instrument outputs (no value typed by hand): the six v0.11.0 instruments (D185, consigned seeds), the v0.12.0 formal model and seal analysis (D187, D189), and the v0.13.0 corrections --- the audited benchmark (the v0.11.0 ``15.3$\times$'' claim is withdrawn as a perimeter artifact, D191) and the corrected seal statistics (the values reported as $\alpha$ in v0.11.0/v0.12.0 were likelihood ratios, D190).}
\label{tab:preuves}
\begin{tabularx}{\textwidth}{@{}p{2.4cm} p{2.7cm} X p{3.6cm}@{}}
\toprule
\textbf{Instrument} & \textbf{Target} & \textbf{Measured} & \textbf{Raw source} \\
\midrule
Mutation testing (16 mutants/block, seed 20261111) & loupe domain (18 sealed blocks) & 230/288 killed (79.9\%); 58 survivors named & \code{preuves\_\allowbreak{}v110/\allowbreak{}mutation.csv} \\
\addlinespace[1.5pt]
Mutation testing (16 mutants/block, seed 20261111) & tls13 domain (28 sealed blocks) & 397/448 killed (88.6\%); 51 survivors named & \code{preuves\_\allowbreak{}v110/\allowbreak{}mutation.csv} \\
\addlinespace[1.5pt]
Mutation testing (16 mutants/block, seed 20261111) & sys domain (6 sealed blocks) & 67/96 killed (69.8\%); 29 survivors named & \code{preuves\_\allowbreak{}v110/\allowbreak{}mutation.csv} \\
\addlinespace[1.5pt]
Mutation testing (same campaign) & cryptographic families (sha-256, hmac/hkdf, x25519, aes-gcm, md5) & 100\% killed (224/224 mutants, 0 survivors) & \code{preuves\_\allowbreak{}v110/\allowbreak{}mutation.csv} \\
\addlinespace[1.5pt]
Coverage-guided fuzzing (seed 20261111) & loupe-01 (untrusted input path) & 72/119 instructions covered (60.5\%); 0 crashes; 0 canary violations; corpus 18 & \code{preuves\_\allowbreak{}v110/\allowbreak{}fuzzage.json} \\
\addlinespace[1.5pt]
Coverage-guided fuzzing (seed 20261112) & http-10 (input parser, sealed 0.4.0) & 358 instructions covered; 0 crashes; 0 named faults; corpus 17 & \code{preuves\_\allowbreak{}v110/\allowbreak{}fuzzage.json} \\
\addlinespace[1.5pt]
Bounded static memory safety (abstract interpretation) & 52 sealed blocks & 2\,098 accesses classified; 1\,181 stack accesses verified, 0 out-of-frame; 263 contract-mediated; 191 constant [rip]; 463 unresolved (named) & \code{preuves\_\allowbreak{}v110/\allowbreak{}surete\_\allowbreak{}memoire.csv} \\
\addlinespace[1.5pt]
Constant-time (emulation dudect, N=60/class, Welch $|t|{>}5$) & tls13-01..06 (sha-256, hmac, hkdf), 13/14 (x25519), 15..19 (aes-gcm), loupe-06 (md5) & all $|t| = 0$ across 6 targets; x25519 executes exactly 1\,238\,959 instructions for every key & \code{preuves\_\allowbreak{}v110/\allowbreak{}temps\_\allowbreak{}constant.csv} \\
\addlinespace[1.5pt]
Oracle independence (fresh seed 20261113) & MD5/SHA-256 tables re-derived (decimal arithmetic, cube roots); RFC 1321 gold vectors spec/harness/hashlib; 200-message os.urandom differentials & 0 divergences across 6 confrontations & \code{preuves\_\allowbreak{}v110/\allowbreak{}independance.csv} \\
\addlinespace[1.5pt]
Comparative benchmark sha256(1,024 B) — AUDITED v0.13.0 (D191) & tls13-01..04 blocks (hand notation, complete operation init+absorb+final) & 75\,983 instructions/op (69.84 instr/round); 765 B code + 256 B table; 27.6 ops/s (emulated) & \code{preuves\_\allowbreak{}v130/\allowbreak{}audit\_\allowbreak{}benchmark.csv} \\
\addlinespace[1.5pt]
Comparative benchmark sha256(1,024 B) — AUDITED v0.13.0 (D191) & gcc -O2 reference (same emulator, same 1,024-byte input) & 69\,132 instructions/op (63.54 instr/round); 1\,226 B code + 256 B table; 30.6 ops/s (emulated) & \code{preuves\_\allowbreak{}v130/\allowbreak{}audit\_\allowbreak{}benchmark.csv} \\
\addlinespace[1.5pt]
Comparative benchmark sha256(1,024 B) — AUDITED v0.13.0 (D191) & gcc -Os reference (same emulator, same input) & 77\,505 instructions/op (71.24 instr/round); 292 B code + 256 B table; 27.4 ops/s (emulated) & \code{preuves\_\allowbreak{}v130/\allowbreak{}audit\_\allowbreak{}benchmark.csv} \\
\addlinespace[1.5pt]
Formal model of CF3 (v0.12.0, D187 — explicit-state exploration) & P1 secret under threshold; P2 roll-back refusal; P3 no destruction on wrong phrase; P4 bounded blast radius & all 4 properties PROVED IN THE MODEL (0 counterexamples); 154 states, 1\,152 transitions, depth 6; adversary at most 2 of 3 elements with epoch replay & \code{preuves\_\allowbreak{}v120/\allowbreak{}modele\_\allowbreak{}formel.csv} \\
\addlinespace[1.5pt]
Seal statistics — CORRECTED v0.13.0 (D190) & rule: seal iff $n$ consecutive identical rebuilds, per-rebuild divergence $q$ & P(seal$|q{=}50\%)$: 3.1\% ($n{=}5$), 0.098\% ($n{=}10$); P(seal$|q{=}1\%$): 95.1\%, 90.4\%; one-sided quantiles 1.86$\sigma$/3.10$\sigma$ & \code{preuves\_\allowbreak{}v130/\allowbreak{}stats\_\allowbreak{}v130.csv} \\
\addlinespace[1.5pt]
RQ1 exact interval (v0.13.0) & 11 of 12 terrain runs surfaced at least one escaped defect & Clopper-Pearson 95\% CI [0.615 — 0.998] (descriptive: sequential iterations of one developer-agent pair) & \code{preuves\_\allowbreak{}v130/\allowbreak{}stats\_\allowbreak{}v130.csv} \\
\addlinespace[1.5pt]
\bottomrule
\end{tabularx}\end{table}


% ------------------------------------------------------------------
\section{Results}\label{sec:results}

\subsection{RQ1: escaped-defect rate}\label{sec:rq1}

Table~\ref{tab:runs} reports the twelve terrain runs of the Loupe case
study. Eleven of the twelve iterations surfaced at least one escaped
defect --- a defect that had passed every laboratory control in force at
the time, including the three Unicorn harnesses (159 controls at three
fixed seeds), the structural oracle (43 then 121 locks) and the seal
guard. The twelfth run (the 0.1.1 artifact, whose card-debt had been paid
at birth) was green end-to-end. The exact Clopper-Pearson 95\% interval
for eleven of twelve is $[0.615, 0.998]$ --- a \emph{descriptive}
binomial summary only: the twelve iterations are sequential deliveries
of one developer--agent pair, not independent trials, so the interval
must not be read as an inference over a population of developers.
Counting named classes across the
campaign: eleven named defect classes over eleven failing iterations,
each with a named mechanism and an archived evidence piece
(\code{preuves-run-loupe-*} in the evidence vault). The showcase ---
the application simply being \emph{used} on the terrain --- was alive with
zero error-reporter events at runs 10, 11 and 12, after six consecutive
deaths.

The earlier milestones show the same shape at smaller scale: the M2
instrumentation campaign took eight iterations and named eight classes
before the hot-patch loop was proven on a real machine (D43); the M3
proxy needed eight terrain deliveries before the first complete run
(D63), and its killer defects (kernel32 shadow-space violations, a
\code{sub rsp,0x7100} without stack probing) were all of classes the
emulator could not see. We treat this as the study's central empirical
finding: \emph{passing a complete laboratory programme does not transfer
to the terrain}, and the protocol's response is not more laboratory
coverage but a terrain evidence regime in which every escape becomes a
named, falsified class --- each of the eleven Loupe classes is now
guarded by a control or an audit born from its own mechanism (e.g.\ the
msx64 ABI audit born at build from the SSE class of run 4).

\subsection{RQ2: token cost per validated change}\label{sec:rq2}

Table~\ref{tab:budgets} shows the instrumented budgets. Locating a block
costs 51 tokens measured against a 64-token design target; editing a
block 719 against 3{,}000; adding a block 1{,}088 against 6{,}000. The
whole generation-1 lineage --- all specifications, drafts and delivered
notations --- totalled 13{,}873 tokens (D34). Under the protocol's
card-fragment indexing, no operation of the update loop ever requires
reading the binary; the card is the only index, and full-card audits are
exceptional, logged operations. These figures predate the Loupe
application and were not re-measured for it; we report them as the
protocol's M1 instrumentation result (a limitation is discussed in
Section~\ref{sec:threats}).

\subsection{RQ3: defect classes invisible to emulation}\label{sec:rq3}

The eleven named classes of the Loupe campaign group into five families;
each was invisible to the emulation harness for a structural reason, not
a coverage gap --- and each corresponds to a guarantee a compiler
normally provides \emph{silently} (a loadable, relocatable image;
ABI-conformant call sequences; conventional error codes; runtime
assumptions baked into the toolchain), which hand notation must
re-derive explicitly, one named lesson at a time:

\begin{enumerate}
\item \emph{Loader and image semantics} (runs 1--3): an ASLR image
without a relocation table (refusal \code{0xc000\allowbreak 007b}); an image outside
the proven envelope of eight prior artifacts (page hole, massive BSS);
an uninitialized \code{.data} session frame (the initializer existed only
in the harness's ideal memory). The harness parses images; it never loads
them.
\item \emph{ABI conventions with hardware consequences} (run 4): msx64
stack arguments 5+ written at the home space instead of
\code{[rsp+0x20]+}, and stack misalignment at call sites
($\mathrm{rsp} \equiv 8 \bmod 16$) --- harmless under emulation, fatal
with a real SSE-using \code{gdi32}. This class gave birth to a static
msx64 ABI audit at build time, which then caught two further alignment
defects no one had seen.
\item \emph{Operating-system loader behaviour} (runs 5, 9, 11): Windows
refusing a write-open on the file of a live mapped section
(\code{gle 32}); import tables walked by the loader at \code{DWORD}
stride while the build wrote \code{QWORD} entries --- the first cause of
all six showcase deaths, found on the artifact's bytes; and the
structural false negative that a library loader never resolves an EXE's
imports (the ``liaison'' probe could never succeed as designed --- the
showcase itself is the proof).
\item \emph{Error-code conventions} (runs 6--7): an \code{hInstance}
read from the PEB instead of \code{GetModule\allowbreak HandleW\allowbreak (NULL)}; negative
return codes read as unsigned 32-bit values by the kit while the piece
delivered them sign-extended.
\item \emph{Harness- and kit-embedded assumptions} (runs 8, 10 and the
pre-run fix of run 11): a probe killed by its own unguarded Python; a
64-bit \code{HMODULE} truncated by \code{ctypes} to a C \code{int}
because \code{argtypes} was not set; a bare \code{ret} neutralizing an
entry point whose return value the loader reads in \code{DllMain}
semantics. These are defects of the \emph{test rig}, named with the same
discipline as defects of the artifact.
\end{enumerate}

The pattern across families: the emulator verifies \emph{logic} well and
\emph{environment} not at all. The protocol's answer is asymmetric ---
add the static audits that can see the environment-sensitive classes
(ABI, image format), and route the rest through terrain evidence with
named-defect accounting, rather than pretending the harness can be
completed.

\subsection{RQ4: hot-patch safety}\label{sec:rq4}

The M2 campaign demonstrated the full loop on a real machine: replacing
the echo block \emph{while executing}, verifying by comparison,
rolling back, re-verifying --- with zero post-patch regression across the
eight instrumented iterations (D43). On RISC-V (M7), where fixed-width
instructions forbid the x86-64 relative trampoline, the patch path uses
a scratch register plus \code{fence.i}. The M7 terrain evidence for this
path was produced on the developer's real machine (WSL2, a real kernel
and real syscalls) executing the RISC-V artifact under the kit's
embarked \code{qemu-riscv64} process emulator --- the instrumental
reserve named in every M7 verdict (D149): the ISA was emulated, never
disguised as native execution. The patch \emph{mechanism} is proven in
that regime (notation, round-trip, the in-traffic switch with byte-exact
transcription, falsification, zero regression; 29/29 recoup on crossed
pieces); the hardware instruction-cache semantics of \code{fence.i} on
physical RISC-V silicon remain \emph{unexecuted} and are explicitly not
claimed. Load-time verification (R10)
checks per-block checksums and signed version monotonicity; the
anti-rollback barrier was exercised at each sealing (reconstruction of a
sealed version is refused, observed at the sealing gesture itself), and
version epochs open strictly monotonically --- the 0.1.1 artifact
$(0,1,1,-1,1)$ opened the new epoch; a pre-release of the sealed 0.1.0
remains refused. Sealed artifacts never changed after their sealing run
(Table~\ref{tab:empreintes}); the 0.1.1 release refreshed a stale
checksum table \emph{at birth} (17 sums, zero bytes of code moved, the
0.1.0 seal intact) --- an explicit debt-repayment pattern the protocol
names and tracks.

% ------------------------------------------------------------------
\subsection{Strengthened evidence (the v0.11.0 programme)}\label{sec:renforce}

The second edition added six laboratory instruments to the evidence base
(Table~\ref{tab:preuves}); all run with fixed, consigned seeds and are
reproducible from the scripts shipped with this edition.

\paragraph{Mutation testing.} 832 single-byte mutants of the sealed blocks
(instructions, operands, offsets --- one byte mutated in a laboratory copy,
the sealed artifact untouched), killed by the existing instruments in
order of cost: the round-trip check, the structural oracle, and --- for
the nineteen pure-compute crypto blocks and loupe-06 --- the semantic
differential against official vectors. Result: 694/832 killed (83.4\%);
the cryptographic families reach 100\%. The 138 survivors concentrate in
the orchestration and UI blocks, where the semantic killer is not wired
at the isolated-mutant level (their semantics are verified by the full
harnesses in composition, not per mutant): this gap is \emph{named}, not
hidden, and is the primary target of the next edition.

\paragraph{Coverage-guided fuzzing.} The two untrusted-input blocks were
attacked under instrumentation: \code{loupe-\allowbreak 01} (the file path, through
the harness's virtual OS with canaries around the path buffer) reached
60.5\,\% block-instruction coverage over 400 iterations with zero crashes
and zero canary violations; the \code{http-10} parser (on the sealed 0.4.0
artifact) reached 358 distinct instructions with zero emulation crashes
and zero named faults, growing a 17-entry corpus. The limits are stated:
coverage is per-block under emulation, and crashes are emulation faults,
not terrain deaths.

\paragraph{Bounded static memory safety.} An abstract interpretation of
all 52 blocks (the extension of the msx64 ABI audit born at run 4):
2{,}098 memory accesses classified --- 1{,}181 stack accesses verified
against the frame bound (zero out-of-frame: the stack-smash class is
absent), 263 accesses mediated by contract-bounded argument pointers, 191
constant \code{[rip]} accesses, 463 unresolved accesses \emph{named}
(volatile-recomputed bases, composed indices). The analysis is linear
(no path enumeration, loop bodies verified once, no alias analysis): a
bounded verification, not a proof.

\paragraph{Constant-time (emulation dudect).} The crypto families pass a
Welch t-test on executed-instruction counts between fixed-secret and
random-secret classes (N=60 each; all $|t| = 0$: the trace is
secret-independent --- e.g.\ x25519 executes exactly 1{,}238{,}959
instructions for every key). The metric is emulation-structural:
\emph{real-hardware wall-clock timing remains unmeasured} and is
explicitly marked as not proven.

\paragraph{Comparative benchmarks --- audited and corrected (v0.13.0).} The
v0.11.0 edition reported a ``15.3$\times$ instruction advantage'' for
the hand-written SHA-256 blocks over a compiled reference. The audit
required by the present edition found that figure to be a \emph{measurement
artifact}: the instrument read the emulation counter after the hash
pipeline returned, but the counter resets at every block invocation ---
the published 4{,}524 instructions were the \emph{finalisation stage
alone}, compared against the reference's \emph{complete} operation (the
implausible 4.2 instructions per SHA-256 round was the tell; a scalar
round needs an order of magnitude more). Re-measured stage by stage,
digest verified against \code{hashlib} before counting and the 17
compression invocations confirmed mechanically, the complete operation
costs 75{,}983 instructions for the blocks against 69{,}132 for
gcc~14.2 -O2 and 77{,}505 for -Os, in the same emulator on the same
input: the hand notation sits \emph{between} the two compiler levels
($-9\%$ in executed instructions against -O2, $+2\%$ against -Os), is
1.45$\times$ more compact than the -O2 function (1{,}021 vs 1{,}482
bytes including the shared K table) but less compact than -Os (548
bytes), and neither camp uses any ISA extension beyond the x86-64
baseline --- no AVX, no SHA-NI (the -O2 reference does use baseline SSE2
moves for buffer setup, which the pure-GPR notation does not; the gap
is reported unadjusted and strict GPR-to-GPR parity is not claimed).
The 15.3$\times$ claim is \emph{withdrawn}; the corrected statement is
that the WYSIWYG notation is competitive with a modern compiler on its
own turf --- the price of direct auditability is nearly nil, and the
spectacular premium does not exist.

\paragraph{Oracle independence.} The oracle was re-generated by routes
independent of the code generator's: the MD5/SHA tables re-derived
through decimal arithmetic, the RFC 1321 gold vectors confronted
spec$\leftrightarrow$\allowbreak harness$\leftrightarrow$\allowbreak hashlib, the differentials
re-run with a new seed and the system generator (\code{os.urandom}):
zero divergences across all confrontations. A second \emph{agent} remains
to be tried (the replication protocol is proposed in the next edition).

% ------------------------------------------------------------------
\section{Discussion}\label{sec:discussion}

\paragraph{What the data supports.} A single developer with an LLM agent
delivered, in 17 days, a cross-platform native codebase with a pure
TLS~1.3 stack and a real Windows application, under a discipline that
made every claim checkable: every defect named, every fix diffed, every
seal fingerprinted, every decision append-only. The escaped-defect data
(RQ1/RQ3) is, to our knowledge, the first quantified account of
\emph{where} agent-written native code fails when laboratory verification
is strong but environmental fidelity is absent: not in logic, but in
loader semantics, ABI conventions and rig assumptions. Agent-failure
quantification exists at source level in repository-scale
settings~\cite{yang2024sweagent}, and the verified-compilation line
quantifies the translator's own guarantees~\cite{leroycompcert}; the
claim is scoped to what is new here --- the sealed-notation level, under
signed-intent governance.

\paragraph{What the data does not support.} The study cannot claim that
the protocol is \emph{more} productive or \emph{more} correct than a
compiler-based workflow: there is no baseline, one case, one agent
family, one developer. The token budgets were measured once, early, on a
small lineage. The oracle is derived from the same layer I it checks, so
specification errors are shared, not caught. Memory safety is enforced
only to the depth of the ABI audits and the fuzzing campaigns;
constant-time is verified \emph{structurally} (secret-independent
instruction traces in emulation) --- real-hardware timing remains
unmeasured. These limits are stated, not smoothed over; the companion
programmes of the two subsequent editions addressed several of them
directly (mutation testing, bounded memory safety and structural
constant-time in v0.11.0; the formal model of the vault chain in
v0.12.0), and the benchmark audit of the present edition is itself an
instance of the method applied to its own claims.

\paragraph{On the ``no compiler'' position.} We phrase the protocol's
position precisely: TRIADE eliminates the \emph{optimizing and opaque}
translation chain by having the agent write the assembly notation itself,
with a round-trip check making the notation WYSIWYG
(what-you-write-is-what-executes) at every build. It does not eliminate
the assembler, the operating system, or the platform ABI --- those are
called through their documented conventions, like any native binary.
Similarly, the sovereignty programme's claim about expropriation is
narrow: the \emph{evolution} of stolen code is blocked (signed layer-I
intent, sealed inheritance, offline keys, dead-man's-switch Shamir
escrow); the \emph{execution} of a stolen copy is not, and the
specification says so.

% ------------------------------------------------------------------
\subsection{Formal model of CF3, structured threat model, and
replication protocol (the v0.12.0 programme)}\label{sec:formel}

\paragraph{A bounded formal model of the CF3 vault.} The sealed-vault
chain (CF3: three-element lock, two-act cycle, epoch anti-rollback,
invocation-based sentinel, destruction on proven falsification) was
modelled as an explicit-state transition system and model-checked
exhaustively (154 states, 1{,}152 transitions, depth 6, three epochs,
adversary interleavings included). The cryptographic primitives are
symbolic (perfect); the adversary holds \emph{at most two} of the three
key elements (all subsets explored), can replay any past
container-plus-trap pair, and can falsify either side. Four claims were
verified, each reported with its status and hypotheses:

\begin{itemize}
\item \textbf{P1 (secret below threshold).} With $\le 2$ of the three
elements, the adversary never reaches the learned-secret state:
\emph{proved in the model} (0 counterexamples). The KEK derivation
binds all three elements (HKDF over $K_{user}\,\|\,S_{lockure}\,\|\,
S_{piege}$ with the epoch as context).
\item \textbf{P2 (rollback refusal).} No transition ever decreases the
epoch; replaying a past container-plus-trap pair is refused by the
crossed-epoch check and the highest-seen mirror: \emph{proved in the
model} (structural monotonicity, mechanically re-checked over every
transition).
\item \textbf{P3 (no destruction on wrong passphrase).} Wrong-passphrase
attempts lead only to the benign retry branch (increasing delay); no
reachable state combines wrong-phrase attempts with local destruction:
\emph{proved in the model} (0 counterexamples).
\item \textbf{P4 (blast radius bounded to one epoch).} Proven
falsification destroys only the local container and draft (with a
report); the trap media is the \emph{witness}, never the accused, and
the CF2 fragments coexist: \emph{proved in the model} (0 violations).
\end{itemize}

The limits are named: the checker is an explicit-state explorer
(Tamarin/ProVerif absent from the laboratory --- the deviation is
consigned); side channels, implementation flaws (RAM discipline, GC),
three-element compromise and quantum adversaries are \emph{out of
model} (hypotheses A1--A5 in the report). ``Proved in the model'' means
exactly that.

\paragraph{Threat model restructured in two planes.} The single threat
table of the specification is restructured along the protocol's own
founding distinction --- obfuscation protects intellectual property,
hardening protects execution --- into two planes:

\begin{itemize}
\item \textbf{Plane IP (intellectual property).} Reverse engineering
(time bought, eroding with LLM decompilers), expropriation (the
evolution of stolen code is blocked --- signed layer-I intent, sealed
inheritance, offline keys, Shamir dead-man's switch --- its execution
is not), redistribution (per-build fingerprints, canary trap).
Confidentiality of the card dominates; loss of confidentiality is
remediable (rotation, variants).
\item \textbf{Plane execution safety.} Card poisoning/corruption
(signature, cross-lock, AEAD), evil-maid on container or trap media
(epoch cross-checks, AEAD tags, the undervivable \code{verif} field),
container-plus-trap rollback (three epoch witnesses), altered binary
(per-block checksums, signed manifest at load), non-conforming
environment (profile checks, fail-closed), agent failure (R2/R5,
oracle from I, budgets). Integrity and availability dominate;
destruction is survivable (CF2 fragments, the trap media).
\end{itemize}

The \emph{intra-Security order} left implicit by the project hierarchy
(D19: Security $>$ Performance $>$ Optimization $>$ Token cost) is now
explicit as project policy: within Security, the IP plane ranks
confidentiality first (remediable losses), while the execution plane
ranks integrity and availability first (irremediable losses); when the
two planes collide, the execution plane wins --- a broken vault protects
nothing. The D81 tension (the control passphrase stored in clear for
survival) and the D183 counter-proposal (encrypted escrow, Shamir
5-of-3) are consigned with the quantified loss/exposure analysis of the
v0.10.0 edition: the expected annual loss under D81 is
$\approx 4.0 \times 10^{-2}\,C_{exposure}$ against
$\approx 1.0 \times 10^{-3}$ under the escrow --- a factor $\approx 40$
--- and D81 \emph{remains} the active decision (the ceremony belongs to
the developer).

\paragraph{The D81 consequence, stated in the model (v0.13.0).} The
honest accounting of the current decision adds a row to the threat
model: the control passphrase lives in clear at the repository (the
survival-over-confidentiality arbitration, D81), and the CF3 container
header carries $S_{serrure}$ --- the machine half of the key material ---
in clear by design. Two of the three key-derivation elements are
therefore \emph{trivially available} to anyone with repository access;
under the threshold adversary of the formal model two elements yield
nothing, but the third --- the trapped medium --- is the only real
confidentiality barrier left, and the integrity guarantees (epoch
witnesses, destruction on proven falsification) carry the protective
load. Two gestures answer this row: the v0.13.0 edition \emph{prepares}
the D183 encrypted-escrow ceremony kit (Shamir 5-of-3 split,
cold-reconstitution proof, media-separation plan, withdrawal checklist)
at the laboratory, \emph{without applying it} --- closing the passphrase
exposure belongs to the developer; and the laboratory's own tests run on
temporary clones with test phrases only (the D105 visible-test boundary:
the real container is never opened by laboratory tests, and test states
are visibly marked as such).

\paragraph{Crystallisation: the statistical model --- corrected (v0.13.0).}
The rule (ten consecutive byte-identical rebuilds before sealing) seals
a build whose per-rebuild divergence probability is $q$ with probability
$(1-q)^n$ (independent rebuilds, the standing hypothesis). Its
operational characteristics: at $q = 50\%$ --- a badly
non-deterministic build --- $n = 5$ would still seal with probability
$3.1\%$ and the historical $n = 10$ with $0.098\%$, a $32\times$
strictness increase, the equivalent of moving from $1.86\sigma$ to
$3.10\sigma$ one-sided; at $q = 1\%$ --- a healthy build --- the rule
seals with probability $95.1\%$ ($n{=}5$) and $90.4\%$ ($n{=}10$). The
values ``$\alpha = 3.29\times10^{-2}$'' and ``$\alpha = 1.08\times
10^{-3}$'' reported in the v0.11.0/v0.12.0 editions were likelihood
ratios $(0.5/0.99)^n$ mislabelled as error probabilities, and their
``sigma'' conversions ($\sqrt{-2\ln p}$) are replaced here by the exact
one-sided normal quantiles. The threshold \emph{remains} ten (project
policy, never changed silently); the corrected analysis is consigned as
the justification and the documented path to any future relaxation.

\paragraph{The replication protocol.} N=1 remains the declared limit;
this edition ships the protocol that would lift it: three independent
axes (a second developer running one full milestone; a second real
application in a different domain; a port to a second non-primary
profile), a T/X coding grid for every defect (five terrain classes T1--T5
--- loader/image, hardware ABI, loader semantics, system conventions,
environment --- against three laboratory classes X1--X3), double
annotation by two independent coders, Cohen's $\kappa$ with the
acceptance threshold $\kappa \ge 0.61$ (substantial agreement), and
pre-registered acceptance criteria. The protocol is proposed, not
executed: the arbitration belongs to the developer.

% ------------------------------------------------------------------
\section{Threats to validity}\label{sec:threats}

\emph{Internal.} The oracle is circular with respect to layer I
(both derive from the same specifications); a specification error passes
the oracle by construction --- the v0.11.0 independence runs (distinct
derivation routes, system generator, spec$\leftrightarrow$\allowbreak harness
$\leftrightarrow$\allowbreak hashlib tripartite confrontation) found zero
divergences, but a second \emph{agent} has not yet regenerated the
oracle. The crystallisation rule (ten identical rebuilds) checks
determinism, not correctness. The mutation-testing survivors (138/832,
concentrated in orchestration/UI blocks without an isolated-mutant
semantic killer) are a named, measured gap --- and the v0.13.0 audit of
our own v0.11.0 benchmark, which caught a perimeter artifact that
inflated a claim fifteen-fold, is itself evidence of this threat class:
instrument defects contaminate results symmetrically, and were caught
here only because the edition required the audit. \emph{External.} N=1: one
project, one developer, one agent family; no baseline comparison;
generalisation is a hypothesis for future work. \emph{Construct.} The
escaped-defect count depends on what counts as a ``named'' class --- we
counted conservatively (a class is counted once even if it killed several
runs); the token budgets measure the M1-era card, not the final
10-domain card, and measure \emph{reads}, not agent reasoning overhead.
\emph{Conclusion.} The terrain is one real Windows machine and one
developer's usage of the application; showcase liveness (0 error-reporter
events) is a weak liveness signal, not a correctness proof. The RISC-V terrain
evidence (M7) was produced under a user-mode emulator with a real
kernel; the \code{fence.i} hardware instruction-cache semantics are
unvalidated on silicon. Memory safety
is \emph{boundedly} verified (2{,}098 accesses classified, zero
stack-smash, 463 named unresolved --- linear analysis, no alias
analysis): a bounded result, not a proof. Constant-time is verified
\emph{structurally} (secret-independent instruction traces); real-hardware
wall-clock behaviour, cache-timing and power side channels remain
UNMEASURED and are marked as not proven.

% ------------------------------------------------------------------
\section{Related work}\label{sec:related}

\paragraph{Design by contract and specification-first development.}
TRIADE's layer I applies design-by-contract ideas~\cite{meyer1992} at the
machine-code level: one specification per block, pre/post-conditions,
named failure cases. The novelty is not the contract style but the
binding of contracts to \emph{byte ranges} through the signed card, and
the rule that no notation byte exists before its specification.

\paragraph{N-version and property-based testing.} The differential
oracles and scenario oracles follow the property-based testing
tradition~\cite{claessen2000} and the classic redundancy arguments of
N-version programming~\cite{avizienis1985}: independent instruments (three
concurrent harnesses, oracles regenerated from I, never exported) must
agree. TRIADE extends the independence requirement to the \emph{terrain}
(a machine the laboratory never controls) and to the \emph{rig itself}
(kit defects are named with the same discipline as artifact defects).

\paragraph{Certification discipline.} The milestone/exit-criteria
structure and the append-only, amend-or-invalidate decision log borrow
from airworthiness process standards such as
DO-178C~\cite{do178c}: objectives with independence, traceability, and
configuration control --- here applied to agent-written native code, with
every criterion measured rather than attested.

\paragraph{Supply-chain integrity.} Per-build fingerprints, signed
manifests, sealed tags and key ceremonies align TRIADE with
reproducible builds \cite{lamb2022}, in-toto's layout
attestations~\cite{intoto2019}, SLSA's build levels~\cite{slsa},
survivable-key-compromise update systems (TUF)~\cite{tuf2010} and
automated signing at scale (Sigstore)~\cite{sigstore2022}; the offline
root-key ceremony follows the discipline of DNSSEC key
ceremonies~\cite{rfc4033}. TRIADE's addition is applying these controls
to \emph{fragments of a development state} (the card) rather than to
released artifacts only, and making the sealed tag a fact of the card
while the artifact's bytes stay frozen. Secret sharing of the master
key~\cite{shamir1979} is used both for escrow and for the dead-man's
switch.

\paragraph{Live patching.} Hot patching of native code has a long
lineage, from dynamic software updating~\cite{hicks2002} and
general-purpose DSU systems such as Kitsune~\cite{hayden2012kitsune} and
POLUS~\cite{polusDSU} to industrial kpatch-style
trampolines~\cite{kpatch}. TRIADE's contribution is the
governance wrapper: a patch is a surgical, diffed, card-tracked revision
of blocks at stable addresses, with crystallisation, rollback and
cross-ISA rules (x86-64 trampoline vs.\ RISC-V scratch+\code{fence.i}),
validated by oracle replay and, since M9, requiring a terrain run
recouped on pieces before any seal.

\paragraph{Repository-scale context for LLM agents.} The token-budget
programme (fragmented card, one fragment per operation, no full-binary
reads) addresses the measured context limits of LLM agents: performance
degradation in long contexts~\cite{liu2024} and the need for
repository-scale retrieval and navigation~\cite{zhang2023, jimenez2024}.
The card is, in effect, a \emph{signed repo-map for machine code}: a
compact index that lets an agent operate on a large binary without ever
loading it.

\paragraph{Learned and verified compilation.} The compiler side of the
argument has moved in two directions: \emph{learned} optimisation, with
LLM-based compiler foundation models trained on optimisation
traces~\cite{cummins2024llmcompiler}, and \emph{verified} compilation,
where the translator is not removed but proven, as in
CompCert~\cite{leroycompcert}. TRIADE takes the opposite door: instead
of making the compiler smarter or proving it correct end to end, it
removes the optimising translator from the agent's path and re-derives
its guarantees explicitly --- the WYSIWYG round-trip at every build,
at a measured token cost. The comparison is not adversarial: CompCert
answers ``can the translation be trusted?'', TRIADE answers ``what does
the agent learn when it must write, and therefore own, the bytes?''

\paragraph{Agent-driven software engineering.} SWE-agent~\cite{yang2024sweagent}
quantifies what interface design does for repository-level agent
performance --- at source-code level, in hosted repositories. The RQ
data of the present study is the sealed-notation counterpart: what
breaks when the same class of agent operates on native code under
signed-intent governance. The two layers are complementary; to our
knowledge no quantified account exists between them.

\paragraph{Simulation-to-reality transfer.} The oracle/terrain
epistemology is a deliberate instance of the reality-gap tradition:
evolutionary robotics learned that controllers evolved in simulation
fail on real hardware unless the gap is itself modelled~\cite{jakobi1995reality};
modern sim-to-real work randomises the simulator to close the
gap~\cite{peng2018simtoreal}, and hardware-in-the-loop validation
interposes real components between the two~\cite{bacic2005hil}. TRIADE's
answer is asymmetric: it does not attempt to close the gap by improving
the emulator, but makes the gap itself the measured object --- every
escape is a named, falsified class with an archived evidence piece
(RQ1), and the RQ3 taxonomy is the resulting map of the gap.

\paragraph{Positioning.} None of the cited primitives is novel. The
claimed contribution is the integration --- contract-first notation at
the byte level, signed-intent governance, oracle/terrain epistemology
with named-defect accounting --- evaluated end-to-end by building a real
application with an LLM agent and reporting where it failed.

% ------------------------------------------------------------------
\section{Conclusion}\label{sec:conclusion}

We reported the first academic edition of TRIADE, a protocol under which
an LLM agent writes native code by hand under signed-intent governance,
and a longitudinal single-case evaluation of its eleven-milestone
validation programme. The data, drawn entirely from the protocol's own
append-only records, quantifies the cost (51/719/1{,}088 tokens per
locate/edit/add), the escape dynamics (eleven named defect classes over
twelve field iterations, all environmental), and the patch safety (zero
post-patch regression, armed anti-rollback) of the approach. The Loupe
application --- a real Windows hex editor born entirely from the method,
sealed twice, and alive on its user's machine --- is the existence proof
that the pipeline can carry a real deliverable. The second edition strengthened the
evidence: 83.4\% mutation kill with the cryptographic families at
100\%, zero-crash coverage-guided fuzzing of the untrusted-input blocks,
bounded memory-safety verification with zero out-of-frame accesses,
secret-independent instruction traces on the crypto families, an audited
instruction-count parity with a compiled SHA-256 reference (within ten
percent of gcc -O2 and ahead of -Os --- the v0.11.0 ``15.3$\times$''
claim withdrawn by the v0.13.0 audit as a measurement artifact), and
zero divergences in oracle-independent regeneration. What remains open
is now precisely measured: the 138 mutation survivors in orchestration
blocks, the emulation-only scope of the constant-time result, the
unmeasured real-hardware timing (including the \code{fence.i} patch
path, whose hardware instruction-cache semantics are unvalidated on
RISC-V silicon), the missing second-agent replication --- and, on the
governance side, the D183 escrow ceremony, prepared and waiting for the
developer's gesture.

% ------------------------------------------------------------------
\section*{Declarations}

\paragraph{Ethics and dual-use.} The artifacts developed under TRIADE
(Loupe, the proxy, the supervisor) are the project's own tools, tested on
the developer's machines; the protocol includes no stealth or
EDR-evasion mechanism, and its specification places furtivity explicitly
out of scope. One counter-measure class is documented honestly in both
documents (the v0.13.0 arbitration, D190): a \emph{light, assumed
anti-debug check by elapsed-time measurement} --- abnormally slow
invocations are treated as dynamic instrumentation --- whose reach is
deterrence, not resistance, and which is consigned as an assumed residual
in the specification's threat table. The work is a dual-use-respectable
engineering study: the same discipline that hardens an artifact against
tampering also makes it auditable. The repository is published under the
Business Source License 1.1 (non-production use until the change date,
then GPL~v3).

\paragraph{Data availability.} All primary data is public in the project
repository (append-only journal, decision log, evidence vault indexes,
extraction script, generated tables). A versioned archival snapshot
(Zenodo DOI) will be minted at submission time and referenced here; the
two submission archives of this edition (\code{TRIADE\_soumission\_v0.13.0.zip}
and \code{TRIADE\_complet\_v0.13.0.zip}) contain respectively the
manuscript with its data tables, and the complete sources with raw
extracts and fingerprints.

\paragraph{Use of AI in this work.} An LLM agent (Z.ai GLM family)
executed the laboratory work of the project under the developer's
instruction, and drafted editorial material under the developer's
direction; all results reported were produced by the protocol's own
instrumentation, and the primary sources (journal, decision log,
evidence vault) were consigned session by session as the work happened.
The agent is not an author. The developer takes full responsibility for
the content of this manuscript.

\paragraph{Acknowledgements.} The author thanks the Z.ai GLM agent
(``the laboratory'') for eleven milestones of disciplined execution, and
the TRIADE decision log for keeping everyone honest.

% ------------------------------------------------------------------
\bibliography{refs}

\end{document}
